% Bagian: sets-functions-relations
% Bab: relations
% Subbagian: reflections
%
\documentclass[../../../include/open-logic-section]{subfiles}

\begin{document}

\olfileid[id]{sfr}{rel}{ref}
\olsection{Refleksi Filosofis}

Dalam \olref[set]{sec}, kita mendefinisikan relasi sebagai himpunan-himpunan
tertentu. Kita perlu berhenti sejenak dan mengajukan sebuah pertanyaan
filosofis singkat: apa yang sebenarnya \emph{dilakukan} oleh definisi semacam
itu? Sangat meragukan bahwa kita ingin mengatakan bahwa kita telah
\emph{menemukan} fakta identitas metafisis tertentu; misalnya, bahwa relasi
urutan pada $\Nat$ \emph{ternyata} merupakan himpunan $R=
\Setabs{\tuple{n,m}}{n, m \in \Nat\text{ dan } n < m}$ yang kita definisikan
dalam \olref[set]{sec}. Berikut tiga alasannya.

Pertama: dalam \olref[set][pai]{wienerkuratowski}, kita mendefinisikan
$\tuple{a, b} = \{\{a\}, \{a, b\}\}$. Sebagai gantinya, tinjaulah definisi
$\lVert a, b\rVert = \{\{b\}, \{a, b\}\} = \tuple{b,a}$. Jika $a \neq b$,
maka $\tuple{a, b} \neq \lVert a,b\rVert$. Namun, kita sama saja dapat
memakai $\lVert a,b\rVert$, bukan $\tuple{a,b}$, sebagai definisi pasangan
terurut. Kedua definisi itu akan bekerja sama baiknya. Jadi, sekarang kita
memiliki dua kandidat yang sama baiknya untuk ``menjadi'' relasi urutan pada
bilangan asli, yaitu:
\begin{align*}
		R &= \Setabs{\tuple{n,m}}{n, m \in \Nat \text{ dan }n < m}\\
		S &= \Setabs{\lVert n,m\rVert}{n, m \in \Nat \text{ dan }n < m}.
\end{align*}
Karena $R \neq S$, berdasarkan ekstensionalitas jelaslah bahwa keduanya tidak
mungkin \emph{sama-sama} identik dengan relasi urutan pada~$\Nat$. Namun,
menyatakan bahwa secara faktual $R$, bukan $S$ (atau sebaliknya),
\emph{merupakan} relasi urutan akan sepenuhnya sewenang-wenang dan karena itu
agak memalukan. (Ini adalah contoh yang sangat sederhana dari argumen
menentang reduksionisme teori himpunan yang dipopulerkan Benacerraf dalam
\citeyear{Benacerraf1965}. Kita akan meninjaunya kembali beberapa kali.)

Kedua: jika kita beranggapan bahwa \emph{setiap} relasi harus diidentifikasi
dengan suatu himpunan, maka relasi keanggotaan himpunan itu sendiri, $\in$,
harus merupakan suatu himpunan tertentu. Himpunan itu haruslah
$\Setabs{\tuple{x,y}}{x \in y}$. Akan tetapi, apakah himpunan ini ada? Mengingat
Paradoks Russell, keberadaan himpunan semacam itu bukanlah klaim sepele.
Faktanya, \oliflabeldef{cumul:::part}{teori himpunan yang kita kembangkan dalam
\olref[cumul][][]{part} akan \emph{menyangkal} keberadaan himpunan ini.
\footnote{Melompat ke depan, berikut alasannya. Untuk reductio, andaikan $I =
\Setabs{\tuple{x,y}}{x \in y}$ ada. Maka $\bigcup \bigcup I$ adalah himpunan
semesta, bertentangan dengan
\olref[sfr][z][sep]{thm:NoUniversalSet}.}}{teori himpunan dapat dikembangkan
secara ketat sebagai teori aksiomatik, dan teori tersebut memang akan
menyangkal keberadaan himpunan ini.}
Jadi, sekalipun beberapa relasi dapat diperlakukan sebagai himpunan, relasi
keanggotaan himpunan harus menjadi kasus khusus.

Ketiga: ketika kita ``mengidentifikasi'' relasi dengan himpunan, kita
memperbolehkan diri menulis $Rxy$ sebagai ganti
$\tuple{x,y} \in R$. Hal ini tidak bermasalah asalkan relasi keanggotaan,
``$\in$'', diperlakukan \emph{sebagai} predikat. Namun, jika kita beranggapan
bahwa ``$\in$'' menyatakan suatu jenis himpunan tertentu, maka ungkapan
``$\tuple{x,y} \in R$'' hanya terdiri atas tiga term tunggal yang menyatakan
himpunan: ``$\tuple{x,y}$'', ``$\in$'', dan ``$R$''. Daftar nama semacam itu
tidak lebih mampu menyatakan sebuah proposisi daripada rangkaian tanpa makna:
``cangkir tempat pena meja''. Sekali lagi, sekalipun beberapa relasi dapat
diperlakukan sebagai himpunan, relasi keanggotaan himpunan harus menjadi kasus
khusus. (Ini menggabungkan versi sederhana paradoks konsep \emph{kuda} Frege
dengan keberatan terkenal Wittgenstein terhadap Russell.)

Lantas, apa kesimpulannya bagi kita? Sebenarnya tidak ada yang \emph{salah}
dengan pernyataan bahwa relasi pada bilangan adalah himpunan. Kita hanya harus
memahami semangat yang melandasi pernyataan tersebut. Kita tidak sedang
menyatakan suatu fakta identitas metafisis. Kita semata-mata mencatat bahwa,
dalam konteks tertentu, kita dapat (dan akan) \emph{memperlakukan} relasi
(tertentu) sebagai himpunan tertentu.

\end{document}
